% Part: normal-modal-logic
% Chapter: completeness
% Section: modalities-ccs

\documentclass[../../../include/open-logic-section]{subfiles}

\begin{document}

\olfileid{nml}{com}{mod}

\olsection{المؤثرات الجهية والمجموعات المتسقة التامة}

\begin{explain}
قد اخترنا للنموذج~$\mModel{M^\Sigma}$ عوالم هي المجموعات
$\Sigma$-المتسقة التامة~${\OLId{greek-variable}{Delta}}$ في منطق جهوي نظامي~$\Sigma$.
وبقي أن نعيّن علاقة إمكان الوصول~${\OLId{latin-looped}{R}}^\Sigma$ بين هذه «العوالم».
وضابط الاختيار، في هذه العلاقة وفي التعيين~${\OLId{latin-looped}{V}}^\Sigma$، أن يتحقق
$\mSat{{\OLId{latin-looped}{M}}^\Sigma}{!A}[{\OLId{greek-variable}{Delta}}]$ إذا وفقط إذا كان~$!A \in {\OLId{greek-variable}{Delta}}$.
فلننظر كيف يهدي هذا الضابط إلى التعريف.

متى عُرّفت علاقة إمكان الوصول، ضمن تعريف الصدق عند عالم أن
\iftag{prvBox}{$\mSat{{\OLId{latin-looped}{M}}^\Sigma}{\Box!A}[{\OLId{greek-variable}{Delta}}]$ إذا وفقط إذا كان
  $\mSat{{\OLId{latin-looped}{M}}^\Sigma}{!A}[{\OLId{greek-variable}{Delta}}']$ لكل~${\OLId{greek-variable}{Delta}}'$ بحيث
  ${\OLId{latin-looped}{R}}^\Sigma{\OLId{greek-variable}{Delta}}{\OLId{greek-variable}{Delta}}'$}{$\mSat{{\OLId{latin-looped}{M}}^\Sigma}{\Diamond!A}[{\OLId{greek-variable}{Delta}}]$ إذا
  وفقط إذا كان~$\mSat{{\OLId{latin-looped}{M}}^\Sigma}{!A}[{\OLId{greek-variable}{Delta}}']$ لبعض~${\OLId{greek-variable}{Delta}}'$ بحيث
  ${\OLId{latin-looped}{R}}^\Sigma{\OLId{greek-variable}{Delta}}{\OLId{greek-variable}{Delta}}'$}. ويتطلب برهان أن
$\mSat{{\OLId{latin-looped}{M}}^\Sigma}{!A}[{\OLId{greek-variable}{Delta}}]$ إذا وفقط إذا كان~$!A \in {\OLId{greek-variable}{Delta}}$ أن
يصدق ذلك خصوصًا في ال!!{formula}s التي تبدأ بمؤثر جهوي؛ أي إن
\iftag{prvBox}{$\mSat{{\OLId{latin-looped}{M}}^\Sigma}{\Box!A}[{\OLId{greek-variable}{Delta}}]$ إذا وفقط إذا كان
  $\Box !A \in {\OLId{greek-variable}{Delta}}$}{$\mSat{{\OLId{latin-looped}{M}}^\Sigma}{\Diamond!A}[{\OLId{greek-variable}{Delta}}]$ إذا
  وفقط إذا كان~$\Diamond !A \in {\OLId{greek-variable}{Delta}}$}. وبجمع هذا المطلب مع تعريف
الصدق عند عالم لـ\iftag{prvBox}{$\Box !A$}{$\Diamond !A$} نحصل على:
\[
\iftag{prvBox}{%
  \Box!A \in {\OLId{greek-variable}{Delta}} \text{ إذا وفقط إذا كان } !A \in {\OLId{greek-variable}{Delta}}'
  \text{ لكل ${\OLId{greek-variable}{Delta}}'$ يحقق } {\OLId{latin-looped}{R}}^\Sigma{\OLId{greek-variable}{Delta}}{\OLId{greek-variable}{Delta}}'}{%
  \Diamond!A \in {\OLId{greek-variable}{Delta}} \text{ إذا وفقط إذا كان } !A \in {\OLId{greek-variable}{Delta}}'
  \text{ لبعض } {\OLId{greek-variable}{Delta}}' \text{ يحقق } {\OLId{latin-looped}{R}}^\Sigma{\OLId{greek-variable}{Delta}}{\OLId{greek-variable}{Delta}}'}
\]
\iftag{prvBox}{تأمل الاتجاه من اليسار إلى اليمين: فهو يقول إنه إذا كان
  $\Box !A \in {\OLId{greek-variable}{Delta}}$، كان~$!A \in {\OLId{greek-variable}{Delta}}'$ لأي~$!A$ وأي~${\OLId{greek-variable}{Delta}}'$
  يحقق~${\OLId{latin-looped}{R}}^\Sigma{\OLId{greek-variable}{Delta}}{\OLId{greek-variable}{Delta}}'$. فإذا اشترطنا أن
  ${\OLId{latin-looped}{R}}^\Sigma{\OLId{greek-variable}{Delta}}{\OLId{greek-variable}{Delta}}'$ إذا وفقط إذا كان~$!A \in {\OLId{greek-variable}{Delta}}'$ لكل
  $\Box!A \in {\OLId{greek-variable}{Delta}}$، تحقق ذلك. ويمكننا أن نكتب الشرط الواقع يمين
  ``إذا وفقط إذا'' بإيجاز أكبر هكذا:
  $\Setabs{!A}{\Box!A \in {\OLId{greek-variable}{Delta}}} \subseteq {\OLId{greek-variable}{Delta}}'$.}{تأمل الاتجاه
  من اليمين إلى اليسار: فهو يقول إنه إذا كان~$!A \in {\OLId{greek-variable}{Delta}}'$، كان
  $\Diamond!A \in {\OLId{greek-variable}{Delta}}$ لأي~$!A$ وأي~${\OLId{greek-variable}{Delta}}'$ يحقق
  ${\OLId{latin-looped}{R}}^\Sigma{\OLId{greek-variable}{Delta}}{\OLId{greek-variable}{Delta}}'$.
  % تصحيح مصدر (0445-I1): التعريف المعلن علاقة تكافؤ لا شرط كفاية وحده.
  فإذا عرّفنا~${\OLId{latin-looped}{R}}^\Sigma{\OLId{greek-variable}{Delta}}{\OLId{greek-variable}{Delta}}'$
  بأنه يتحقق إذا وفقط إذا كان~$\Diamond!A \in {\OLId{greek-variable}{Delta}}$ لكل~$!A \in {\OLId{greek-variable}{Delta}}'$،
  لزم هذا الاتجاه من التعريف.
  ويمكننا أن نكتب الشرط الواقع يمين ``إذا وفقط إذا'' بإيجاز أكبر هكذا:
  $\Setabs{\Diamond!A}{!A \in {\OLId{greek-variable}{Delta}}'} \subseteq {\OLId{greek-variable}{Delta}}$.}

ويبقى السؤال: هل يضمن تعريف~${\OLId{latin-looped}{R}}^\Sigma$ هذا بالفعل أن
\iftag{prvBox}{$\Box !A \in {\OLId{greek-variable}{Delta}}$ إذا وفقط إذا كان
  $\mSat{{\OLId{latin-looped}{M}}^\Sigma}{\Box !A}[{\OLId{greek-variable}{Delta}}]$؟ \iftag{prvDiamond}{وهل يضمن
    أيضًا أن }{}}{}\iftag{prvDiamond}{$\Diamond !A \in {\OLId{greek-variable}{Delta}}$ إذا
  وفقط إذا كان~$\mSat{{\OLId{latin-looped}{M}}^\Sigma}{\Diamond !A}[{\OLId{greek-variable}{Delta}}]$؟}{}
ستثبت النتائج القليلة الآتية ذلك.
\end{explain}

\begin{defn}
  إذا كانت~${\OLId{greek-variable}{Gamma}}$ مجموعة من ال!!{formula}s، فلتكن
  \begin{align*}
    \Box{\OLId{greek-variable}{Gamma}} & = \Setabs{\Box !B}{!B \in {\OLId{greek-variable}{Gamma}}}\\
    \Diamond{\OLId{greek-variable}{Gamma}} & = \Setabs{\Diamond !B}{!B \in {\OLId{greek-variable}{Gamma}}}\\
    \intertext{ولتكن}
    \Box^{-{\OLMathNumeral{1}}}{\OLId{greek-variable}{Gamma}} & = \Setabs{!B}{\Box !B \in {\OLId{greek-variable}{Gamma}}}\\
    \Diamond^{-{\OLMathNumeral{1}}}{\OLId{greek-variable}{Gamma}} & = \Setabs{!B}{\Diamond !B \in {\OLId{greek-variable}{Gamma}}}\\
  \end{align*}
\end{defn}

فالمقصود بـ~$\Box{\OLId{greek-variable}{Gamma}}$ أخذ~${\OLId{greek-variable}{Gamma}}$ ووضع~$\Box$ أمام كل
!!{formula} من~${\OLId{greek-variable}{Gamma}}$. وأما~$\Box^{-{\OLMathNumeral{1}}}{\OLId{greek-variable}{Gamma}}$ فنأخذ لتكوينها
ال!!{formula}s التي تبدأ بـ~$\Box$ في~${\OLId{greek-variable}{Gamma}}$، ثم نحذف~$\Box$ الأولى
من كل واحدة. وإنما أدخلنا هذا الاصطلاح اختصارًا للترميز،
لا لأن في هذا التعريف وحده نتيجة جديدة مهمة.

لاحظ أن~$\Box\Box^{-{\OLMathNumeral{1}}}{\OLId{greek-variable}{Gamma}} \subseteq {\OLId{greek-variable}{Gamma}}$:
\[
\Box\Box^{-{\OLMathNumeral{1}}}{\OLId{greek-variable}{Gamma}} = \Setabs{\Box!B}{\Box!B \in {\OLId{greek-variable}{Gamma}}}
\]
أي إنها ليست إلا مجموعة جميع ال!!{formula}s في~${\OLId{greek-variable}{Gamma}}$ التي تبدأ
بـ~$\Box$.

\begin{lem}\ollabel{lem:box1}
  إذا كان~${\OLId{greek-variable}{Gamma}} \Proves[\Sigma] !A$، فإن
  $\Box{\OLId{greek-variable}{Gamma}} \Proves[\Sigma] \Box !A$.
\end{lem}

\begin{proof}
  من~${\OLId{greek-variable}{Gamma}} \Proves[\Sigma] !A$ نختار مقدمات متناهية
  $!B_{\OLMathNumeral{1}}$، \dots،~$!B_{\OLId{latin-ordinary}{k}} \in {\OLId{greek-variable}{Gamma}}$ تحقق
  $\Sigma \Proves !B_{\OLMathNumeral{1}} \lif (!B_{\OLMathNumeral{2}} \lif \cdots (!B_{\OLId{latin-ordinary}{k}} \lif !A)\cdots)$.
  ولأن~$\Sigma$ نظامية، يصح استعمال قاعدة~\RK{}، فنحصل على
  $\Sigma \Proves \Box!B_{\OLMathNumeral{1}} \lif (\Box!B_{\OLMathNumeral{2}} \lif \cdots
  (\Box!B_{\OLId{latin-ordinary}{k}} \lif \Box!A)\cdots)$.
  والمقدمات الجديدة $\Box!B_{\OLMathNumeral{1}}$، \dots،~$\Box!B_{\OLId{latin-ordinary}{k}} \in \Box{\OLId{greek-variable}{Gamma}}$
  بحكم التعريف. فهذا الاشتقاق بعينه شاهد على
  $\Box{\OLId{greek-variable}{Gamma}} \Proves[\Sigma] \Box !A$.
\end{proof}

\begin{lem}\ollabel{lem:box2}
  إذا كان~$\Box^{-{\OLMathNumeral{1}}}{\OLId{greek-variable}{Gamma}} \Proves[\Sigma] !A$، فإن
  ${\OLId{greek-variable}{Gamma}} \Proves[\Sigma] \Box!A$.
\end{lem}

\begin{proof}
  لنفرض $\Box^{-{\OLMathNumeral{1}}}{\OLId{greek-variable}{Gamma}} \Proves[\Sigma] !A$.
  بتطبيق \olref{lem:box1} يلزم
  $\Box\Box^{-{\OLMathNumeral{1}}}{\OLId{greek-variable}{Gamma}} \Proves[\Sigma] \Box!A$.
  وقد علمنا أن $\Box\Box^{-{\OLMathNumeral{1}}}{\OLId{greek-variable}{Gamma}} \subseteq {\OLId{greek-variable}{Gamma}}$؛
  فيجوز توسيع المقدمات، وبالرتابة نحصل على
  ${\OLId{greek-variable}{Gamma}} \Proves[\Sigma] \Box!A$.
\end{proof}

\iftag{prvBox}{%
% Box is primitive

\begin{prop}\ollabel{prop:box}
  إذا كانت~${\OLId{greek-variable}{Gamma}}$ $\Sigma$-متسقة تامة، فإن~$\Box !A \in {\OLId{greek-variable}{Gamma}}$
  إذا وفقط إذا كان، لكل مجموعة $\Sigma$-متسقة تامة~${\OLId{greek-variable}{Delta}}$ تحقق
  $\Box^{-{\OLMathNumeral{1}}}{\OLId{greek-variable}{Gamma}} \subseteq {\OLId{greek-variable}{Delta}}$، لدينا~$!A \in {\OLId{greek-variable}{Delta}}$.
\end{prop}

\begin{proof}
  افترض أن~${\OLId{greek-variable}{Gamma}}$ $\Sigma$-متسقة تامة. الاتجاه ``فقط إذا'' يسير:
  افترض أن~$\Box !A \in {\OLId{greek-variable}{Gamma}}$ وأن~$\Box^{-{\OLMathNumeral{1}}}{\OLId{greek-variable}{Gamma}} \subseteq
  {\OLId{greek-variable}{Delta}}$. لما كان~$\Box!A \in {\OLId{greek-variable}{Gamma}}$، كان
  $!A \in \Box^{-{\OLMathNumeral{1}}}{\OLId{greek-variable}{Gamma}} \subseteq {\OLId{greek-variable}{Delta}}$، ومن ثم~$!A \in {\OLId{greek-variable}{Delta}}$.

  وأما اتجاه «إذا» فنأخذ مقابله. لنفرض $\Box!A \notin {\OLId{greek-variable}{Gamma}}$.
  بما أن~${\OLId{greek-variable}{Gamma}}$ $\Sigma$-متسقة تامة، فهي مغلقة استنتاجيًا؛
  ومن ثم ${\OLId{greek-variable}{Gamma}} \Proves/[\Sigma] \Box !A$.
  ويقتضي مقابل~\olref{lem:box2} أن
  $\Box^{-{\OLMathNumeral{1}}}{\OLId{greek-variable}{Gamma}} \Proves/[\Sigma] !A$.
  وبحسب \olref[prf][con]{prop:consistencyfacts}\olref[prf][con]{prop:consistencyfacts-b}
  تكون~$\Box^{-{\OLMathNumeral{1}}}{\OLId{greek-variable}{Gamma}} \cup \{\lnot!A\}$ $\Sigma$-متسقة.
  فنوسّعها بمبرهنة ليندنباوم إلى مجموعة $\Sigma$-متسقة تامة~${\OLId{greek-variable}{Delta}}$
  تحقق $\Box^{-{\OLMathNumeral{1}}}{\OLId{greek-variable}{Gamma}} \cup \{\lnot!A\} \subseteq {\OLId{greek-variable}{Delta}}$.
  ولما كانت مشتملة على النفي ومتسقة، كان $!A \notin {\OLId{greek-variable}{Delta}}$،
  فوجدنا المجموعة التي تنقض الشرط الكلي.
\end{proof}
}{%
% Box is not primitive, only Diamond is

\begin{prop}\ollabel{prop:diamond}
  إذا كانت~${\OLId{greek-variable}{Gamma}}$ $\Sigma$-متسقة تامة، فإن~$\Diamond !A \in {\OLId{greek-variable}{Gamma}}$
  إذا وفقط إذا وجدت مجموعة $\Sigma$-متسقة تامة~${\OLId{greek-variable}{Delta}}$ تحقق
  $\Diamond{\OLId{greek-variable}{Delta}} \subseteq {\OLId{greek-variable}{Gamma}}$ و~$!A \in {\OLId{greek-variable}{Delta}}$.
\end{prop}

\begin{proof}
  افترض أن~${\OLId{greek-variable}{Gamma}}$ $\Sigma$-متسقة تامة. الاتجاه من اليمين إلى اليسار
  يسير: لتكن~${\OLId{greek-variable}{Delta}}$ مجموعة $\Sigma$-متسقة تامة تحقق
  $\Diamond{\OLId{greek-variable}{Delta}} \subseteq {\OLId{greek-variable}{Gamma}}$ و~$!A \in {\OLId{greek-variable}{Delta}}$. عندئذ
  $\Diamond!A \in \Diamond{\OLId{greek-variable}{Delta}} \subseteq {\OLId{greek-variable}{Gamma}}$، أي
  $\Diamond!A \in {\OLId{greek-variable}{Gamma}}$.

  أما الاتجاه من اليسار إلى اليمين، فافترض أن
  $\Diamond !A \in {\OLId{greek-variable}{Gamma}}$. علينا أن نبين وجود مجموعة $\Sigma$-متسقة
  تامة~${\OLId{greek-variable}{Delta}}$ بحيث~$!A \in {\OLId{greek-variable}{Delta}}$ و
  $\Diamond{\OLId{greek-variable}{Delta}} \subseteq {\OLId{greek-variable}{Gamma}}$. لما كان~$\Diamond !A \in {\OLId{greek-variable}{Gamma}}$
  وكانت~${\OLId{greek-variable}{Gamma}}$ $\Sigma$-متسقة، كان~$\Box\lnot !A \notin {\OLId{greek-variable}{Gamma}}$.
  ولما كانت~${\OLId{greek-variable}{Gamma}}$ $\Sigma$-متسقة تامة، كانت مغلقة استنتاجيًا، ومن
  ثم~${\OLId{greek-variable}{Gamma}} \Proves/[\Sigma] \Box\lnot !A$. وبحسب~\olref{lem:box2}،
  $\Box^{-{\OLMathNumeral{1}}}{\OLId{greek-variable}{Gamma}} \Proves/[\Sigma] \lnot !A$. وبحسب
  \olref[prf][con]{prop:consistencyfacts}\olref[prf][con]{prop:consistencyfacts-b}،
  تكون~$\Box^{-{\OLMathNumeral{1}}}{\OLId{greek-variable}{Gamma}} \cup \{!A\}$ $\Sigma$-متسقة. وبمبرهنة
  ليندنباوم، توجد مجموعة $\Sigma$-متسقة تامة~${\OLId{greek-variable}{Delta}}$ بحيث
  $\Box^{-{\OLMathNumeral{1}}}{\OLId{greek-variable}{Gamma}} \cup \{!A\} \subseteq {\OLId{greek-variable}{Delta}}$. ومن الواضح أن
  $!A \in {\OLId{greek-variable}{Delta}}$. ولإثبات~$\Diamond{\OLId{greek-variable}{Delta}} \subseteq {\OLId{greek-variable}{Gamma}}$، افترض
  خلافه: يوجد~$!B \in {\OLId{greek-variable}{Delta}}$ بحيث~$\Diamond !B \notin {\OLId{greek-variable}{Gamma}}$.
  ولما كانت~${\OLId{greek-variable}{Gamma}}$ $\Sigma$-متسقة تامة، كان
  $\Box\lnot !B \in {\OLId{greek-variable}{Gamma}}$. لكن عندئذ أيضًا
  $\lnot !B \in \Box^{-{\OLMathNumeral{1}}}{\OLId{greek-variable}{Gamma}} \subseteq {\OLId{greek-variable}{Delta}}$، وهذا يجعل
  ${\OLId{greek-variable}{Delta}}$ غير متسقة.
\end{proof}
}

\begin{lem}\ollabel{lem:box-iff-diamond}
  افترض أن~${\OLId{greek-variable}{Gamma}}$ و${\OLId{greek-variable}{Delta}}$ مجموعتان $\Sigma$-متسقتان تامتان. عندئذ
  $\Box^{-{\OLMathNumeral{1}}}{\OLId{greek-variable}{Gamma}} \subseteq {\OLId{greek-variable}{Delta}}$ إذا وفقط إذا كان
  $\Diamond{\OLId{greek-variable}{Delta}} \subseteq {\OLId{greek-variable}{Gamma}}$.
\end{lem}

\begin{proof}
  لإثبات اتجاه «فقط إذا»، نفرض~$\Box^{-{\OLMathNumeral{1}}}{\OLId{greek-variable}{Gamma}} \subseteq {\OLId{greek-variable}{Delta}}$،
  ونأخذ~$\Diamond!A \in \Diamond{\OLId{greek-variable}{Delta}}$، أي~$!A \in {\OLId{greek-variable}{Delta}}$.
  المطلوب $\Diamond!A \in {\OLId{greek-variable}{Gamma}}$؛ ويكفي له~$\Box\lnot!A \notin {\OLId{greek-variable}{Gamma}}$،
  إذ تلزم عندئذ~$\lnot\Box\lnot !A \in {\OLId{greek-variable}{Gamma}}$ بالقصوية.
  ولو كان~$\Box\lnot!A \in {\OLId{greek-variable}{Gamma}}$ لكان
  $\lnot!A \in {\OLId{greek-variable}{Delta}}$ بفرض الاحتواء، فيبطل اتساق~${\OLId{greek-variable}{Delta}}$،
  لأن~$!A \in {\OLId{greek-variable}{Delta}}$. فثبت~$\Box\lnot!A \notin {\OLId{greek-variable}{Gamma}}$ والمطلوب.

  ولإثبات اتجاه «إذا»، نفرض~$\Diamond{\OLId{greek-variable}{Delta}} \subseteq {\OLId{greek-variable}{Gamma}}$
  ونستعمل المقابل: نأخذ~$!A \notin {\OLId{greek-variable}{Delta}}$ ونثبت~$\Box!A \notin {\OLId{greek-variable}{Gamma}}$.
  فمن~$!A \notin {\OLId{greek-variable}{Delta}}$ يلزم~$\lnot!A \in {\OLId{greek-variable}{Delta}}$ بالقصوية؛
  وبفرض الاحتواء يلزم~$\Diamond\lnot!A \in {\OLId{greek-variable}{Gamma}}$.
  غير أن~$\Diamond\lnot!A$ تكافئ~$\lnot\Box !A$ في كل منطق جهوي نظامي.
  فتكون هذه الأخيرة في~${\OLId{greek-variable}{Gamma}}$ بالإغلاق، ويمنع الاتساق خلافها:
  $\Box!A \notin{\OLId{greek-variable}{Gamma}}$.
\end{proof}

\iftag{prvBox}{%
  \iftag{prvDiamond}{%
% Box and Diamond are both primitive; prove
% prop:diamond using lem:box-iff-diamond

\begin{prop}\ollabel{prop:diamond}
  إذا كانت~${\OLId{greek-variable}{Gamma}}$ $\Sigma$-متسقة تامة، فإن~$\Diamond !A \in {\OLId{greek-variable}{Gamma}}$
  إذا وفقط إذا وجدت مجموعة $\Sigma$-متسقة تامة~${\OLId{greek-variable}{Delta}}$ تحقق
  $\Diamond{\OLId{greek-variable}{Delta}} \subseteq {\OLId{greek-variable}{Gamma}}$ و~$!A \in {\OLId{greek-variable}{Delta}}$.
\end{prop}

\begin{proof}
افترض أن~${\OLId{greek-variable}{Gamma}}$ $\Sigma$-متسقة تامة. لدينا
$\Diamond!A \in {\OLId{greek-variable}{Gamma}}$ إذا وفقط إذا كان
$\lnot\Box\lnot !A \in {\OLId{greek-variable}{Gamma}}$، بحسب~\Dual{} والإغلاق.
ولدينا~$\lnot\Box\lnot !A \in {\OLId{greek-variable}{Gamma}}$ إذا وفقط إذا كان
$\Box\lnot !A \notin {\OLId{greek-variable}{Gamma}}$، بحسب
\olref[ccs]{prop:ccs-properties}\olref[ccs]{prop:ccs-lnot}، لأن
${\OLId{greek-variable}{Gamma}}$ $\Sigma$-متسقة تامة. وبحسب~\olref{prop:box}، يكون
$\Box\lnot !A \notin {\OLId{greek-variable}{Gamma}}$ إذا وفقط إذا وجدت مجموعة $\Sigma$-متسقة
تامة~${\OLId{greek-variable}{Delta}}$ تحقق~$\Box^{-{\OLMathNumeral{1}}}{\OLId{greek-variable}{Gamma}} \subseteq {\OLId{greek-variable}{Delta}}$ و
$\lnot!A \notin {\OLId{greek-variable}{Delta}}$. تأمل الآن أي~${\OLId{greek-variable}{Delta}}$ كهذه. بحسب
\olref{lem:box-iff-diamond}، يكون~$\Box^{-{\OLMathNumeral{1}}}{\OLId{greek-variable}{Gamma}} \subseteq {\OLId{greek-variable}{Delta}}$
إذا وفقط إذا كان~$\Diamond{\OLId{greek-variable}{Delta}} \subseteq {\OLId{greek-variable}{Gamma}}$. كذلك، يكون
$\lnot !A \notin {\OLId{greek-variable}{Delta}}$ إذا وفقط إذا كان~$!A \in {\OLId{greek-variable}{Delta}}$، بحسب
\olref[ccs]{prop:ccs-properties}\olref[ccs]{prop:ccs-lnot}. وهكذا،
$\Diamond!A \in {\OLId{greek-variable}{Gamma}}$ إذا وفقط إذا وجدت مجموعة $\Sigma$-متسقة
تامة~${\OLId{greek-variable}{Delta}}$ تحقق~$\Diamond{\OLId{greek-variable}{Delta}} \subseteq {\OLId{greek-variable}{Gamma}}$ و
$!A \in {\OLId{greek-variable}{Delta}}$.
\end{proof}

}{}}{}

\begin{prob}
  بيّن أنه إذا كانت~${\OLId{greek-variable}{Gamma}}$ $\Sigma$-متسقة تامة، فإن
  \iftag{prvBox}{$\Diamond !A \in {\OLId{greek-variable}{Gamma}}$ إذا وفقط إذا وجدت مجموعة
    $\Sigma$-متسقة تامة~${\OLId{greek-variable}{Delta}}$ بحيث
    $\Box^{-{\OLMathNumeral{1}}}{\OLId{greek-variable}{Gamma}} \subseteq {\OLId{greek-variable}{Delta}}$ و~$!A \in {\OLId{greek-variable}{Delta}}$.}{$\Box !A
    \in {\OLId{greek-variable}{Gamma}}$ إذا وفقط إذا كان، لكل مجموعة $\Sigma$-متسقة تامة
    ${\OLId{greek-variable}{Delta}}$ تحقق~$\Box^{-{\OLMathNumeral{1}}}{\OLId{greek-variable}{Gamma}} \subseteq {\OLId{greek-variable}{Delta}}$، لدينا
    $!A \in {\OLId{greek-variable}{Delta}}$.} \emph{افعل ذلك من دون استعمال
    \olref[nml][com][mod]{lem:box-iff-diamond}}.
\end{prob}

\end{document}
